\begin{abstract}
Concept unlearning in text-to-image diffusion models aims to suppress a target concept, e.g., \texttt{horse}, while preserving the model's ability to generate semantically related but distinct content, like \texttt{donkey}. Yet existing methods either leak under indirect prompts or visibly degrade remaining concepts. 
{\em We show that these failure modes arise naturally from overlapping concept representations.} Our main contribution is demonstrating that the observed trade-offs stem from the geometry of concept representations, rather than from weaknesses in any particular algorithm.
We support this claim through both theoretical analysis and empirical observations. By formalizing concepts as activation-space regions, we show that overlap between a target and other concepts lower-bounds the unavoidable degradation on those other concepts when the target is erased. 
Our analysis suggests that strong erasure and preservation become fundamentally coupled when concepts occupy overlapping activation regions, with the trade-off scaling linearly in the degree of overlap. We verify this trade-off on seven unlearning methods spanning fine-tuning, adversarially-robust fine-tuning, and inference-time interventions. Averaged across target concepts, STEREO almost completely suppresses the target concept under indirect prompts but cuts the model's ability to generate semantically related concepts by more than 75\%. On the other hand, sparse inference-time methods (SAeUron, SEOT) better preserve utility but leave substantial target leakage. The trade-off is steep for concepts whose internal representations heavily overlap with their evaluated semantic neighborhoods, (e.g., \texttt{horse} with \texttt{cat}, \texttt{dog}, and \texttt{bear}) and milder for \texttt{castle}, which is comparatively less entangled within our concept pool and probe layer in activation space. Across targets, damage to related concepts scales monotonically with our activation-overlap measure. These results suggest that perfect unlearning is the wrong target for entangled concepts. The field should evaluate on the Pareto frontier our theorem establishes; current benchmarks, which decouple unlearning accuracy from utility preservation, hide this trade-off and need to be revised.
\end{abstract}
%{\em We prove this is unavoidable.} Our main contribution is showing that the failure modes arise from the geometry of concept representations, rather than from weaknesses in any particular algorithm. 
% Thus the perfect erasure and utility preservation cannot coexist whenever the target overlaps with other concepts,
% \varun{don't give a count here; will bias reviewers negatively}
% \ali{that is probably not true in any dataset, so maybe mention or clarify that for which dataset castle is well-separated?} 
% Our main contribution is an impossibility theorem showing that the failure modes arise from the geometry of concept representations, not weaknesses of any particular algorithm. \ali{maybe something like this would be better: Our main contribution is showing that the failure modes arise from the geometry of concept representations, rather than from weaknesses in any particular algorithm. We support this claim through both theoretical analysis and empirical observations.}

% Unlike classical machine unlearning, where the forget and retain sets are typically disjoint, concept unlearning operates on a representation where related concepts share substantial structure:
% we show that semantically related concepts share overlapping activation regions in diffusion models, so the assumption that erasure and neighbor retention can be optimized independently is fundamentally violated.
% This entanglement produces two coupled failure modes: \emph{concept leakage}, where an erased concept resurfaces from prompts that omit its name but supply correlated context (e.g., generating a horse from ``a jockey at the racetrack''), and \emph{collateral forgetting}, where pushing erasure harder degrades neighboring concepts the user wanted to keep.

% We formalize this coupling and prove that no unlearning method can simultaneously achieve perfect erasure (suppressing the target on \emph{all} prompts) and perfect preservation of its semantic neighbors, with the unavoidable trade-off scaling linearly with the degree of representational overlap between the target and its neighbors.
% any method that suppresses the target concept's generation probability below $\varepsilon$ on \emph{all} prompts must displace the activations of its semantic neighbors by at least $\kappa(1-\varepsilon)/L - \delta$ where $L$ is a Lipschitz constant of the generation pipeline and $\delta$ is the inter-concept activation distance. Setting $\varepsilon = 0$ yields a strict impossibility: perfect erasure and perfect neighbor preservation cannot coexist whenever $\kappa > 0$. 
% achieving $\epsilon$-robust erasure \varun{what does the eps refer to?} must incur a non-zero error \varun{what does error mean here?} in preserving neighbors, which is a fundamental trade-off between robustness and retention. 


